\documentclass{article}
\usepackage{hyperref}

\title{论文提纲}
\author{}
\date{\today}

\begin{document}

\maketitle

\tableofcontents

\section{绪论}
\subsection{研究背景与意义}
\subsection{国内外研究现状}
\subsection{研究内容与技术路线}

\section{机场牵引车辆系统建模与问题描述}
\subsection{机场牵引作业场景抽象}
\subsection{坐标系与状态变量定义}
\subsection{车辆运动学模型建立}
\subsubsection{非完整约束条件}
\subsubsection{单轨等效模型}
\subsubsection{双桥转向角速度表达式}
\subsection{转弯半径与曲率模型}
\subsection{典型转向模式分析}
\subsubsection{前桥单独转向模式}
\subsubsection{前后桥反向转向模式}
\subsubsection{前后桥同向转向模式}
\subsection{速度、转角与曲率约束建模}
\subsubsection{速度约束}
\subsubsection{转向角约束}
\subsubsection{曲率约束}
\subsection{车辆动力学简化与驱动关系}
\subsection{状态方程统一表达}
\subsection{仿真实现中的参数映射关系}
\subsection{本章小结}

\section{基于飞机辅助感知的自动对接识别方法}
\subsection{感知任务分析与总体方案设计}
\subsection{飞机辅助感知目标与识别信息设计}
\subsection{感知系统组成与传感器配置}
\subsubsection{视觉识别模块}
\subsubsection{近距测量模块}
\subsection{目标身份确认模型}
\subsection{相对位姿感知模型}
\subsection{近距离态势感知与安全判定}
\subsection{感知流程设计}
\subsection{本章小结}

\section{基于相对位姿误差的自动接近与对接控制方法}
\subsection{控制任务描述与总体思路}
\subsection{状态变量、控制输入与误差定义}
\subsubsection{车辆状态与控制输入}
\subsubsection{感知误差接口}
\subsection{对接过程分析与分层控制架构}
\subsubsection{对接过程分阶段描述}
\subsubsection{分层控制架构}
\subsection{误差动力学模型建立}
\subsection{局部参考轨迹生成方法}
\subsubsection{预对接点与对接点构造}
\subsubsection{三次多项式局部轨迹}
\subsection{参考速度调度策略}
\subsection{误差反馈控制律设计}
\subsubsection{纵向速度控制律}
\subsubsection{角速度控制律}
\subsubsection{对接过程中的控制行为分析}
\subsection{四轮转向映射与约束处理}
\subsubsection{曲率计算}
\subsubsection{双桥反向转向模式}
\subsubsection{前桥单转向模式}
\subsubsection{速度与转角}
\subsection{终端停车判定与模式切换}
\subsection{控制算法主流程}
\subsection{本章小结}

\section{仿真系统构建与实现}
\subsection{系统总体架构与软硬件环境}
\subsection{牵引车模型与传感器建模}
\subsection{任务与模式管理}
\subsection{ROS与Gazebo接口实现}
\subsection{本章小结}

\end{document}